\documentclass[12pt]{article}

% revise margins
\setlength{\headheight}{0.0in}
\setlength{\topmargin}{0.0in}
\setlength{\headsep}{0.0in}
\setlength{\textheight}{9.0in}
\setlength{\footskip}{0.0in}
\setlength{\oddsidemargin}{0.0in}
\setlength{\evensidemargin}{0.0in}
\setlength{\textwidth}{6.5in}

\pagestyle{empty}
\usepackage{times}
\usepackage{amsmath}

\begin{document}

\noindent \textbf{Biostat 140.668} \hfill Problem set 2

\bigskip 

\noindent \textbf{Hidden Markov models}

\begin{enumerate}

\item Consider a single intercross individual derived from two inbred
mouse strains.  Consider $M$ ordered markers, and let $G_i \in \{AA,
AB, BB\}$ denote the (phase-unknown) genotype of the individual at
marker $i$.  Show that the $G_i$ form a Markov chain.

\item Fill in the details in the Viterbi algorithm to calculate
$$ \hat{g} = \arg \max_{g_1, \dots, g_M} \Pr(G_1 = g_1, \dots, G_M =
g_M \ | \ \boldsymbol{O})$$

\end{enumerate}




\end{document}
